\section{Гомоморфизмы простых расширений}

Задача: понять, какие бывают $F$-гомоморфизмы из простых расширений.

\begin{stmt}
    Пусть $F$ --- поле, $F(\delta)$ --- простое расширение, $L/F$ --- другое расширение. Тогда:

    \begin{enumerate}
        \item Если $\delta$ --- трансцендентный элемент, то для любого $F$-гомоморфизма $\varphi\colon F(\delta)
        \to L$ элемент $\varphi(\delta)$ трансцендентен, и $\varphi$ полностью определяется значением
        $\varphi(\delta)$. Имеется биекция:
        \[
            \{F\text{-гомоморфизмы } F(\delta) \to L\} \longleftrightarrow \{\text{трансцендентные элементы }
            L \text{ над } F\}.
        \]

        \item Если $\delta$ --- алгебраический элемент с минимальным многочленом $f(x) \in F[x]$, то для любого
        $F$-гомоморфизма $\varphi\colon F[\delta] \to L$ элемент $\varphi(\delta)$ --- корень $f(x)$.
        Имеется биекция:
        \[
            \{F\text{-гомоморфизмы } F[\delta] \to L\} \longleftrightarrow \{\text{корни } f(x) \text{ в }
            L\}.
        \]
    \end{enumerate}

    \begin{proof}
        Рассмотрим гомоморфизм подстановки $\theta\colon F[x] \to L$, $p(x) \mapsto p(\delta)$.

        \textbf{Случай 1: $\delta$ трансцендентный.} Тогда $\ker\theta = 0$, поэтому $\theta$ --- инъекция и
        $F[x] \cong F[\delta]$. Любой $F$-гомоморфизм $\varphi\colon F(\delta) \to L$ определяется значением
        $\varphi(\delta)$, и $\varphi(\delta)$ должен быть трансцендентным (иначе ядро было бы ненулевым)
        . Обратно, для любого трансцендентного $d \in L$ можно определить $\varphi(\delta) = d$ и продолжить на
        $F(\delta) = \{p(\delta)/q(\delta)\}$ как $\varphi(p(\delta)/q(\delta)) = p(d)/q(d)$.

        \textbf{Случай 2: $\delta$ алгебраический.} Тогда $\ker\theta = (f(x))$ --- главный идеал, порождённый минимальным многочленом
        $f$. Так как $f$ неприводим, $(f(x))$ --- максимальный идеал, и
        \[
            F[\delta] \cong \frac{F[x]}{(f(x))} \text{ --- поле.}
        \]

        Для $F$-гомоморфизма $\varphi\colon F[\delta] \to L$: так как $f(\delta) = 0$, применяя $\varphi
        $, получаем $0 = \varphi(f(\delta)) = f(\varphi(\delta))$, то есть $\varphi(\delta)$ --- корень
        $f(x)$.

        Обратно, пусть $d$ --- корень $f(x)$ в $L$. Построим $F$-гомоморфизм $F[x] \to L$, $x \mapsto d
        $. Его ядро содержит $(f(x))$ (так как $f(d) = 0$). По теореме о гомоморфизме спускается до $F[x]
        /(f(x)) \to L$, что даёт $\varphi\colon F[\delta] \to L$.
    \end{proof}
\end{stmt}


\section{Поля разложения}

\begin{stmt}
    Пусть $F$ --- поле, $f(x) \in F[x]$. Тогда существует поле разложения $E_f$ многочлена $f$ такое,
    что $[E_f : F] \leq (\deg f)!$.

    \begin{proof}
        Пусть $p(x)$ --- неприводимый делитель $f(x)$ степени $n$. Рассмотрим $d_1$ --- корень $p(x)$ в расширении
        $F[d_1] \cong F[x]/(p(x))$. Тогда $[F[d_1] : F] \leq n$.

        В поле $F[d_1]$ многочлен $f(x)$ имеет корень $d_1$. Делим: $f_2(x) = f(x)/(x - d_1)$. Повторяем:
        находим корень $d_2$ многочлена $f_2$ и получаем $F[d_1, d_2]$, затем $d_3$, и так далее.

        Дойдём до $f_{\deg f + 1}$ --- константы. В этот момент $E_f = F[d_1, \dots, d_n]$.

        Считаем степени:
        \[
            [F[d_1] : F] \leq n, \quad [F[d_1, d_2] : F[d_1]] \leq n-1, \quad \dots
        \]
        \[
            [E_f : F] = [F[d_1, \dots, d_n] : F] \leq n \cdot (n-1) \cdot \ldots \cdot 1 = n! \leq (\deg f)
            !. \qedhere
        \]
    \end{proof}
\end{stmt}


\section{Сепарабельность}

\begin{definition}
    Пусть $f(x) \in F[x]$, $d$ --- корень $f(x)$. Говорят, что $d$ --- \emph{кратный корень}, если $f(x)
    = (x - d)^n g(x)$, где $g(d) \neq 0$ и $n \geq 2$.
\end{definition}

\begin{stmt}
    $d$ --- кратный корень $f(x)$ тогда и только тогда, когда $d$ --- корень $f'(x)$.
    \begin{proof}
        Если $f(x) = (x - d)^n g(x)$, то
        \[
            f'(x) = n(x - d)^{n-1} g(x) + (x - d)^n g'(x).
        \]
        При $x = d$: $f'(d) = n \cdot 0^{n-1} \cdot g(d) = 0$ при $n \geq 2$. Если же $n = 1$, то $f'(d)
        = g(d) \neq 0$.
    \end{proof}
\end{stmt}


\begin{stmt}
    Пусть $f(x)$ --- неприводимый многочлен в $F[x]$. Тогда следующие условия эквивалентны:
    \begin{enumerate}
        \item $f$ имеет кратный корень.
        \item $\gcd(f, f') \neq 1$.
        \item Поле $F$ имеет характеристику $p \neq 0$ и $f$ --- многочлен от $x^p$.
        \item Все корни $f$ кратные.
    \end{enumerate}

    \begin{proof}
        $(1 \Rightarrow 2)$. Пусть $d$ --- кратный корень $f(x)$. Тогда $d$ --- корень $f'(x)$. Значит
        $(x - d)$ делит и $f(x)$, и $f'(x)$, откуда $\gcd(f, f') \neq 1$.

        $(2 \Rightarrow 3)$. Так как $f$ неприводим и $\gcd(f, f') \neq 1$, а $\deg f' < \deg f$, единственная возможность ---
        $f' \equiv 0$.

        Если $f(x) = \sum_{i=0}^n a_i x^i$, то $f'(x) = \sum_{i=1}^n i \cdot a_i x^{i-1} = 0$. Значит
        $i \cdot a_i = 0$ для всех $i = 1, \dots, n$.

        Если $\mathrm{char}\, F = 0$, то $a_i = 0$ для всех $i \geq 1$, и $f$ --- константа. Противоречие.
        Значит $\mathrm{char}\, F = p > 0$.

        Тогда $i \cdot a_i = 0$ означает: либо $p \mid i$, либо $a_i = 0$. То есть $f(x) = a_0 + a_p x^p + a_{2p}
        x^{2p} + \dots$, и $f$ зависит только от $x^p$.

        $(3 \Rightarrow 4)$. $f(x) = g(x^p)$. В поле разложения $g(x) = \prod_i (x - \delta_i)^{k_i}$.
        Тогда
        \[
            f(x) = g(x^p) = \prod_i (x^p - \delta_i)^{k_i} = \prod_i (x^p - \beta_i^p)^{k_i} = \prod_i (x - \beta_i)
            ^{pk_i},
        \]
        где $\beta_i^p = \delta_i$ (в ещё большем расширении). Здесь использовано, что $(a + b)^p = a^p + b^p
        $ в характеристике $p$. Все кратности $\geq p$, то есть все корни кратные.

        $(4 \Rightarrow 1)$. Очевидно.
    \end{proof}
\end{stmt}


\begin{cor}
    Пусть $f \neq 0$, $f \in F[x]$. Тогда следующие условия эквивалентны:
    \begin{enumerate}
        \item $\gcd(f, f') = 1$ в $F[x]$.
        \item $f$ не имеет кратных корней (в поле разложения).
    \end{enumerate}
\end{cor}


\begin{definition}
    Многочлен $f$ называется \emph{сепарабельным}, если $\gcd(f, f') = 1$, то есть все корни $f(x)$ простые.

\end{definition}

\begin{definition}
    Элемент $d \in E$ расширения $E/F$ называется \emph{сепарабельным}, если его минимальный многочлен над
    $F$ сепарабелен.
\end{definition}

\begin{definition}
    Расширение $E/F$ называется \emph{сепарабельным}, если $\forall\, d \in E$ элемент $d$ сепарабелен.

\end{definition}

\begin{definition}
    Поле $F$ называется \emph{совершенным}, если $\forall$ неприводимый $f(x) \in F[x]$ является сепарабельным.

\end{definition}

\begin{remark}
    Любое поле характеристики $0$ совершенно. Конечные поля тоже совершенны. Несовершенные поля возникают только в характеристике
    $p > 0$ и обязательно бесконечны.
\end{remark}
